\documentclass[aspectratio=169,12pt]{beamer}

\geometry{paperwidth=200mm,paperheight=112.5mm}
\usepackage{lmodern}
\usepackage{fontspec}
\usepackage{polyglossia}
\setdefaultlanguage{russian}
\setotherlanguage{english}
\setmainfont{Arial}
\setsansfont{Arial}
\newfontfamily\cyrillicfont{Arial}
\newfontfamily\cyrillicfontsf{Arial}
\usefonttheme{professionalfonts}
\usepackage{amsmath,amssymb,booktabs,tabularx}
\usepackage{pgfplots}
\pgfplotsset{compat=1.18}

\definecolor{ink}{HTML}{182C43}
\definecolor{accent}{HTML}{007F82}
\definecolor{muted}{HTML}{5C6875}
\definecolor{dense}{HTML}{B86442}
\definecolor{lightline}{HTML}{DFE5E9}
\setbeamercolor{normal text}{fg=ink,bg=white}
\setbeamercolor{structure}{fg=accent}
\setbeamercolor{frametitle}{fg=ink,bg=white}
\setbeamercolor{title}{fg=ink}
\setbeamercolor{itemize item}{fg=accent}
\setbeamercolor{itemize subitem}{fg=accent}
\setbeamercolor{footline}{fg=muted}
\setbeamerfont{frametitle}{size=\fontsize{20}{23}\selectfont,series=\bfseries}
\setbeamerfont{title}{size=\fontsize{25}{29}\selectfont,series=\bfseries}
\setbeamerfont{footline}{size=\fontsize{7.5}{9}\selectfont}
\setbeamersize{text margin left=9mm,text margin right=9mm}
\setbeamertemplate{navigation symbols}{}
\setbeamertemplate{itemize item}{\small\textbullet}
\setbeamertemplate{itemize subitem}{\small\textbullet}
\setbeamertemplate{frametitle}{\vspace{3mm}\insertframetitle\par\vspace{1mm}}
\setbeamertemplate{footline}{%
  \hfill
  \insertframenumber\,/\,\inserttotalframenumber\hspace{9mm}\vspace{3mm}}
\setlength{\parskip}{2pt}
\AtBeginDocument{%
  \setlength{\abovedisplayskip}{7pt}%
  \setlength{\belowdisplayskip}{7pt}%
  \setlength{\abovedisplayshortskip}{4pt}%
  \setlength{\belowdisplayshortskip}{5pt}}
\newcommand{\pos}[1]{\left[#1\right]_{+}}
\newcommand{\source}[1]{\par\vspace{1.5mm}{\fontsize{7.5}{9}\selectfont\color{muted}#1\par}}
\newcommand{\takeaway}[1]{\par\vspace{2mm}{\color{accent}\bfseries #1\par}}
\newcommand{\paper}{\href{https://arxiv.org/pdf/2601.22766v3}{Santos et al., arXiv:2601.22766v3}}
\newcommand{\mm}{Memory Mosaics}

\begin{document}

% 1
\begin{frame}[plain]
  \vspace{5mm}
  {\usebeamerfont{title}Sparse Attention as\\[1mm]
  Compact Kernel Regression\par}
  \vspace{5mm}
  {\small Saul Santos, Nuno Gonçalves, Daniel C. McNamee,\\
  Marcos Treviso, André F. T. Martins\par}
  \vspace{5mm}
  {\small\phantom{Егор Прокопов}%
    \begin{tikzpicture}[remember picture,overlay]
      \node[anchor=south east,inner sep=0pt,font=\small]
        at ([xshift=-9mm,yshift=9mm]current page.south east) {Егор Прокопов};
    \end{tikzpicture}}
\end{frame}

% 2
\begin{frame}{Главная идея}
  Модель может лучше извлекать нужную информацию, если присваивать достаточно непохожим
  ключам ровно нулевой вес.

  Результат работы внимания - получение взвешенной суммы векторов-значений:
  $$
    y=\sum_{i=1}^{n}p_i v_i,\qquad p_i\geq 0,\qquad \sum_i p_i=1.
  $$

  В обычном softmax-механизме внимания веса определяются нормированными скалярными
  произведениями запроса и ключа. Каждый ключ получает положительный вес.
  Если ключей много, то их суммарный вклад "ненужных" может уменьшать влияние "полезных". 
\end{frame}

% 3
\begin{frame}{Внимание как ядерная регрессия}
  Веса внимания задает функция близости между запросами и ключами - кернел.
  $$
    p_i=\frac{\kappa(q,k_i)}{\sum_j\kappa(q,k_j)},
    \qquad y(q)=\sum_i p_i v_i.
  $$
  
  Кернел $\kappa\geq 0$ определяет, насколько сильно ключ влияет на ответ.

  При $\|q\|=\|k_i\|=1$ сходство и расстояние связаны:
  $$
    s_i=q^{\mathsf T}k_i,
    \qquad d_i^2 = \|q-k_i\|^2 = 2 - 2qk^\top = 2 - 2s_i.
  $$
\end{frame}

% 4
\begin{frame}{Гауссовское ядро даёт softmax-внимание}
  Используем ядро с шириной \(h>0\):
  
  $$
    \kappa_{h}(q,k_i)=\exp\!\left(-\frac{\|q-k_i\|^2}{2h^2}\right).
  $$
  
  Для единичных норм:
  $$
    \kappa_{h}(q,k_i)
       =\underbrace{\exp(-1/h^2)}_{\text{общий множитель}}
        \exp(s_i/h^2).
  $$

  Общий множитель сокращается при нормировании:
  $$
    p_i=\frac{\exp(s_i/h^2)}{\sum_j\exp(s_j/h^2)}
       =\operatorname{softmax}(s/h^2)_i.
  $$
\end{frame}

% 5
\begin{frame}{Компактное ядро обнуляет дальние ключи}
  \begin{columns}[T,onlytextwidth]
    \begin{column}{0.49\textwidth}
      Ядро Епанечникова:
      $$
        \kappa_h(q,k_i)=
        \pos{1-\frac{d_i^2}{h^2}}
      $$

      $[x]_+=\max(0,x)$

      \(h>0\) — радиус отсечения:
      $$
        d_i\geq h\ \Longrightarrow\ p_i=0
      $$

      Близость определяется в пространстве
      ключей. Можно сказать, что позиция не участвует в вычислении, 
      но она участвует опосредованно, через трансформирование векторов посредством 
      позиционного кодирования
    \end{column}
    \begin{column}{0.48\textwidth}
      \centering
      \begin{tikzpicture}
        \begin{axis}[
          width=\linewidth,height=46mm,
          xmin=0,xmax=2.5,ymin=0,ymax=1.08,
          xlabel={\(d/h\)},ylabel={\(\kappa\)},
          tick label style={font=\scriptsize},
          label style={font=\small},
          axis line style={lightline},
          grid=major,grid style={lightline},
          legend style={font=\fontsize{7.5}{9}\selectfont,draw=none,
            at={(0.98,0.97)},anchor=north east},
          samples=100]
          \addplot[dense,very thick,domain=0:2.5]{exp(-x*x/2)};
          \addlegendentry{Гауссовское}
          \addplot[accent,very thick,domain=0:1]{1-x*x};
          \addplot[accent,very thick,domain=1:2.5,forget plot]{0};
          \addlegendentry{Епанечникова}
        \end{axis}
      \end{tikzpicture}
      {\fontsize{8}{10}\selectfont\color{muted}
        Формы ядер без констант нормировки.\par}
    \end{column}
  \end{columns}
  Важно: замена ядра меняет механизм внимания, оно не тождественно softmax-вниманию.
\end{frame}

% 6
\begin{frame}{Разреженное внимание через компакт-ядра: фиксированный радиус через нормированный ReLU}
  Подставим $d_i^2=2-2s_i$ в компактное ядро:
  $$
    w_i=\pos{1-\frac{2-2s_i}{h^2}}
       =\frac{2}{h^2}\pos{s_i-c},
    \qquad c=1-\frac{h^2}{2}.
  $$
  После сокращения общего множителя:
  $$
    p_i=\frac{\pos{s_i-c}}{\sum_j\pos{s_j-c}}.
  $$
  Порог $c$ задаётся заранее выбранным радиусом $h$. 
  Число активных ключей зависит от их сходства с запросом.
\end{frame}

% 7
\begin{frame}{Адаптивный радиус приводит к sparsemax}
  Введём температуру $\gamma>0$ и перепишем те же веса:
  $$
    w_i=\frac{2\gamma}{h^2}
         \pos{\frac{s_i}{\gamma}-\tau},
    \qquad \tau=\frac{1}{\gamma}-\frac{h^2}{2\gamma}.
  $$
  
  Выберем $\tau$ так, чтобы сумма выражений в скобках была единицей:
  $$
    \sum_i\pos{\frac{s_i}{\gamma}-\tau}=1.
  $$

  Тогда после нормирования:
  $$
    p_i=\pos{\frac{s_i}{\gamma}-\tau}
       =\operatorname{sparsemax}(s/\gamma)_i,
    \qquad h^2=2-2\gamma\tau.
  $$
  Таким образом, порог и соответствующий ему радиус подстраиваются под текущие ключи.
\end{frame}

% 8
\begin{frame}{Top-k наиболее похожих ключей}
  Сначала вычисляем скалярное произведение запроса с каждым ключом:
  $$
    s_i=q^{\mathsf T}k_i,\qquad
    S_k=\operatorname{TopK}(s),\qquad 1\leq k\leq n.
  $$
  $S_k$ - индексы $k$ ключей с наибольшими значениями $s_i$.
  При единичных нормах это также $k$ ближайших к запросу ключей.

  Остальным ключам присваиваем нулевой вес.
  
  top-k softmax нормирует веса выбранных ключей:
  $$
    p_i=
    \begin{cases}
      \displaystyle\frac{\exp(s_i/\gamma)}
        {\sum_{j\in S_k}\exp(s_j/\gamma)}, & i\in S_k,\\[2mm]
      0, & i\notin S_k,
    \end{cases}
    \qquad \gamma>0.
  $$
  В uniform top-k все выбранные ключи получают одинаковый вес $1/k$.
\end{frame}

% 9
\begin{frame}{Сравнение Top-k и ReLUmax}
  Top-k выбирает $k$ ключей с наибольшим сходством.
  ReLUmax оставляет ключи, достаточно близкие по сходству к лучшему:
  
  $$
    m = \max_j s_j,\qquad
    w_i = \pos{b+\frac{s_i-m}{h^2}}^r,
  $$
  
  $$
    p_i = \frac{w_i}{\sum_j w_j},\quad b,h,r>0.
  $$

  Ненулевой вес получают ключи, для которых
  
  $$
    s_i > m - bh^2.
  $$
  
  ReLUmax фиксирует допустимое отставание от максимума. 
  Если почти лучших ключей много, то в вычислении участвуют многие ключи.
  Если почти лучших ключей мало, то в вычислении участвует мало ключей (преимущественно - только лучшие).

\end{frame}

% 10
\begin{frame}{Итоги}
  \begin{itemize}
    \setlength{\itemsep}{4mm}
    
          \item Разреженное внимание можно представить
          через компакт-кернел. 
    
          \item Отсечение слабых совпадений помогает
          в ряде задач, но результат зависит от нормировки и самой задачи.
    
          \item Нулевые веса можно использовать для ускорения вычислений, 
          но подходяшего раздела в статье нет

  \end{itemize}
\end{frame}

\end{document}
